\chapter{Background}

\begin{chapterintro}
This chapter gives information about the organisations the researcher worked with as well as background information about unemployment and job readiness (JR) programs.
\end{chapterintro}

\section{Context: Unemployment in South Africa}

The United Nations Sustainable Development Goals names “Decent Work and Economic Growth” as goal number eight, noting that widening inequalities has led to slow job creation, with over 204 million persons unemployed in 2015. There are many factors that cause unemployment: slow job creation, educational inequalities and social instability. In South Africa in particular -- there are many more unskilled labourers than positions available for them, while positions requiring skilled labour often go unfilled. These potential employees, further, face many barriers to employment. Transport costs are exacerbated by geographic inequalities persisting from forced relocations during apartheid. Reliance on an unreliable transport system deepens their difficulties; transport to and from work takes a long time, and train delays make them late for work despite early starts. Youth and foreigners especially struggle with language inequalities; from online job postings, many employers require fluency in English and Afrikaans, rather than isiXhosa, which is the medium of teaching in many township-based schools, or the native languages of African refugees seeking employment in South Africa. Finally, foreigners face difficulties getting visas to work, forcing them to take informal employment at wages below living costs, affecting their performance at work and their ability to retain jobs. At the same time, because of word-of-mouth based recommendations, those that can get and keep a job for a year are often those that stay employed for the rest of their lives.

\section{Job Readiness Programs}

Many non-governmental organisations (NGOs) address unemployment through job readiness (JR) programs, which provide soft-skills training to help their participants improve their English, learn how to obtain and keep jobs, and seek additional training in specific trades. These programs include some computer skills training, often necessary for curriculum vitae preparation and for finding available posts online, but also a valuable skill in general, making learners more employable. In this section, we describe the JR programs of our two partner NGOs.

\subsection{The Zanokhanyo Network (TZN)}

The Zanokhanyo Network (TZN) is a non-profit organisation targeting unemployed individuals. TZN assists to provide employee-sought skills and self-confidence boost amongst other benefits to all kinds of job seekers, both foreign nationals and locals. We focused primarily on TZN’s Wynberg site, which serves people of various ages from all over Cape Town, including many French-speaking refugees. Other courses offered in this site include a basic English language course, computer skills training, waiter/waitressing, and barista courses.

The JR course runs over a 12-day period in which the following modules are delivered:

\begin{enumerate}[label=\textbf{\arabic*.}]
    \item Who am I?
    \item Emotional Healing
    \item Why character matters
    \item Introduction to computers
    \item Stumbling blocks on the journey, but God\ldots
    \item Self-management
    \item Searching for employment
    \item Applying for vacancies
    \item Getting your job
    \item At the workplace
    \item Landmines in the workplace
    \item Beyond Job Readiness
\end{enumerate}

After completing the course, the JR graduates are invited to return any time to make free use of the resource centre, which is an internet-connected computer lab. The JR classes were offered for R60 (~\$5) per student. The job readiness classes are few to enable the students to cover the whole syllabus before rushing to find a job, also to reach as many students as can be reached annually, which is in line with South Africa’s vision 2030 of providing one million learning opportunities through community education and training centres \cite{npc2013}. The more students reached by each job-readiness programme, the quicker the goal will be reached.

\subsection{Afrika Tikkun (AT)}

Like TZN, AT is an international organisation that trains individuals to become ready for employment. AT works in partnership with SEDA and ESKOM. AT runs three core programs suited to different age groups:

\begin{enumerate}[label=\textbf{\arabic*.}]
\item Early Childhood Development program
\item Child and Youth Development program
\item Career Development program
\end{enumerate}

In this research, we focus on the Career Development Program, which is a two-month JR program, which targets 19--35 year old individuals. The NGO trains individuals by providing them with money skills, people skills and work management skills. AT runs their JR classes at Mfuleni, a township nearly 40km from Cape Town. The readiness classes are offered free to the students.

\section{Participatory Design for Development}

Participatory Design is a design approach that attempts to obtain bi-directional knowledge transfer between users and designers (Information Technology experts), for users to have a say and allow the designer to make the possible approaches to solving problems known to the users. Although PD was targeted at office work \cite{simonsen2013}, it has shifted to use in research. PD allows users to be more engaged, “having a say” \cite{simonsen2013} as opposed to being “spoken for” by the researcher \cite{frascara2002} in a user-centred design approach. In PD the user is recognized as an expert \cite{bjorgvinsson2010}.

PD involves in-depth knowledge often from few individuals, as such PD has sometimes been criticized as being small scaled and short-lived \cite{simonsen2013}, mapping only to a small or local context and dying once the project is completed. We acknowledge this criticism as a critical ethical issue for both PD and Co-design; however, we reduce the chance of running into this issue by working with NGOs as partners capable of disseminating information about and mediating the use of artefacts generated during research over different intakes of students.

\section{Co-design for Development}

Co-design is a process of collective creativity throughout the design process \cite{sanders2008} among different kinds of experts \cite{steen2011}.

Like Co-design, Participatory Design (PD) is a design approach that attempts to obtain bi-directional knowledge transfer between users and designers. However, Co-design has a significant difference in that all participants are considered equals in the design process, irrespective of their role \cite{zimmerman2011}. Allowing equity for all participants ensures that the non-IT-experts, in this case job seekers, can have a strong voice in the design process.

The participants in Co-design can be academics, researchers, designers or other professionals; however, all are experiential experts, that is, experts of their own experiences \cite{steen2011} and of their preferences \cite{preece2002}, which are useful for coming up with designs and choosing among multiple designs.

In terms of living labs, there are two “research types”: observation research, where a user is considered as a subject, and participative research, where a user actively contributes in co-creating value. Our methodology, Participatory Design, and our design approach, Co-design, both fall under participative research \cite{pallot2011}.

We consider job seekers as experts in understanding the challenges in job searching and as such should be deeply involved in the research and not just be informants. Engaging users early in the design process allows users to be more likely to identify issues that can be met by technology usage and increases their likelihood to adopt a design solution into their practices \cite{david2013}.

Co-design is very critical to research that is aimed at development because although some values are known with little input directly from participants, not all values are explicit and common to individuals. Co-design provides a setup where users can be challenged by other users or through the performance of certain tasks to express themselves \cite{halloran2017}. Co-design enables several benefits: it fosters joint creativity \cite{steen2011}, induces ownership and sustainability \cite{david2013}, and promotes the generation of inspired ideas through probes \cite{sanders2014}.

In our Co-design cycles, we make use of various tools and techniques such as workshops, interviews, walkthroughs and participant observation. The Co-design methodology promotes new ways of engagement and as such should be flexible to change \cite{sanders2008}. Thus, we make use of the Tell $\leftrightarrow$ Make $\leftrightarrow$ Enact Co-design cycle which is flexible to design phase changes.

We aimed to achieve beneficiary empowerment \cite{ssozi2016}, stakeholder equity \cite{zimmerman2011}, cost optimisation \cite{steen2011} and project sustainability \cite{david2013}, which are all possible when using the Co-design methodology.